\scanpage{204}
\mspage{13}
% archive sheet number: 203

Pour $C_0$ parcourant les cribles \deleted{\uline{couvrants}}\inserted{admissibles} de $C$, \deleted{les $R$} et admettant que la catégorie \inserted{fibrée} \uline{$C$ est loc\textsuperscript{t}} \uline{filtrante}, ce qui implique que l'intersection de deux cribles couvrants est encore couvrant, on voit que les $R_S^{C_0}$ forment une famille filtrante de ss-catégories multigaloisiennes pleines de $R_S$, de réunion $R_S$. Donc on gagne si on sait que les foncteurs $\alpha$ sont des \uline{équivalences} de catégories, car on aura alors une équivalence

\[
R_S\approxeq\mathop{\underrightarrow{\mathrm{Lim}}}\limits_{\substack{C_0\\\text{variable}}}\ \mathop{\underleftarrow{\mathrm{Lim}}}\limits_{\text{sur }C_0}(\unsure{\mathcal E|E}|C_0)
\]

qui \unsure{explicitent|explicite} $R_S$ en termes de $(C,s)$, \unsure{sous réserve|pourvu vérifiée|qui reste} tout au moins que la notion de \uline{crible admissible}

% margin note written sideways
\margin{soit connue. Il en sera ainsi si on impose que pour tt $M\in\mathrm{Ob}\,C_{S'}$, $(S'\in\mathrm{Ob}\,\mathrm{Et}(S))$ $\varphi(M)\longrightarrow e_{S'}$ \uline{est épimorphisme dans $R_{S'}$} ; \uline{c'est ce que nous supposerons par la suite} pour simplifier (Mais ce n'est pas très naturel pour les \unsure{mouvements|monuments} \unsure{alg.} ! On en doit \unsure{pouvoir} \unsure{justifier|particulariser} \unsure{G.R.} !)}

On va montrer que pour que $\alpha$ soit une équivalence (pour $C_0$ crible couvrant) il suffit que $\varphi$ soit \uline{pleinement fidèle}.

\margin{\unsure{Supposons|Supposant} \unsure{tjs} $C$ loc. filtrante}

Comme $\alpha$ est un foncteur exact de catégories multigaloisiennes, il suffit de voir pour la \deleted{\unsure{plei|plus}}fidélité que $\alpha(X)=\emptyset\Longrightarrow X=\emptyset$, ce qui est immédiat grâce au fait $C_0$ \deleted{couvrant}\inserted{admissible} : c'est local pour (Et), on peut supposer donc qu'il existe un $M\in C_{0S}$, \inserted{avec $\varphi(M)\longrightarrow e_S$ épim,} alors comme $r(\varphi(M)\times X)=\emptyset$ on a $\varphi(M)\times X=\emptyset$, OK.

$\alpha$ pleinement fidèle ? Soit $X\in\mathrm{Ob}\,R_S^{C_0}$, et $Z=\alpha(X)=Z_1\amalg Z_2$, montrer que cette

